\chapter{Software Development}
\label{chap:software-development}

\section{Software Development Methodology}
\label{section:software-development-methodology}

This project follows Extreme Programming (XP), an agile 
software development methodology that emphasizes short 
development cycles, continuous feedback, and close 
collaboration between team members and stakeholders. XP 
is suitable for this project because the team is small, 
the development timeline is fixed, and the technical 
requirements involve frequent experimentation and 
iteration, particularly in the adversarial generation 
pipeline.

The following characteristics of this project align 
with XP practices:

\begin{itemize}
    \item \textbf{Small team:} The project is developed 
    by two members, which supports close communication 
    and a shared understanding of all components without 
    requiring heavy coordination overhead.
    \item \textbf{Academic oversight:} Three supervisors 
    provide consistent review and feedback throughout 
    the development process, serving a role similar to 
    the customer representative in XP.
    \item \textbf{Fixed timeline:} The project has a 
    defined development period of approximately four 
    months, which encourages prioritization and 
    incremental delivery of working features.
    \item \textbf{Single codebase:} The entire system 
    is maintained in a single shared repository, 
    supporting continuous integration and consistent 
    versioning across the frontend and backend components.
    \item \textbf{Continuous testing:} Adversarial 
    generation results and system behavior are evaluated 
    regularly throughout development to verify that the 
    pipeline produces the expected outputs.
\end{itemize}

\section{Technology Stack}
\label{section:technology-stack}
<TIP: Describe your technology stack here. See the following example from ThaiProgrammer.org />
\begin{figure}[h]
    \centering
    \includegraphics[width=0.5\textwidth]{chapter_6/Tech-Stack.png}
    \caption{technology stack}
\end{figure}

\section{Coding Standards}
\label{section:coding-standards}
<TIP: Describe your coding standard for this project here. />

\section{Progress Tracking Report}
\label{section:progress-tracking-report}
<TIP: Show that you have been working on this project overtime.
It can be in the form of a burndown chart or a contribution graph from GitHub./>
